% Part: computability
% Chapter: computability-theory
% Section: universal-part-function

\documentclass[../../../include/open-logic-section]{subfiles}

\begin{document}

\olfileid{cmp}{thy}{uni}
\olsection{الدالة الجزئية الشاملة القابلة للحساب}

\begin{thm}
\ollabel{thm:univ-comp}
توجد دالة جزئية شاملة قابلة للحساب~$\fn{Un}(e,x)$. وبعبارة أخرى، توجد
دالة $\fn{Un}(e,x)$ تحقق:
\begin{enumerate}
\item $\fn{Un}(e,x)$ جزئية قابلة للحساب.
\item إذا كانت $f(x)$ أي دالة جزئية قابلة للحساب، وجد عدد طبيعي $e$
بحيث $f(x)\simeq\fn{Un}(e,x)$ لكل~$x$.
\end{enumerate}
\end{thm}

\begin{proof}
لتكن $\fn{Un}(e,x)\simeq U(\umin{s}{T(e,x,s)})$، حيث $U$ و$T$ كما
في مبرهنة كليني للصورة السوية (\olref[nfm]{thm:normal-form}).
\end{proof}

\begin{explain}
ليست هذه إلا طريقة دقيقة للقول إن لدينا تعدادًا فعالًا للدوال الجزئية
القابلة للحساب. والفكرة هي أننا إذا كتبنا $f_e$ للدالة المعرفة بـ
$f_e(x)=\fn{Un}(e,x)$، فإن المتتالية $f_0$ و$f_1$ و$f_2$ و… تضم
جميع الدوال الجزئية القابلة للحساب، مع خاصية إمكان حساب $f_e(x)$
«بانتظام» في $e$ و~$x$. وللتبسيط نستعمل دالة ثنائية شاملة للدوال
الأحادية الحجة، لكن يمكننا بسهولة، بترميز متتاليات الأعداد، تعميم ذلك
على عدد أكبر من الحجج. فلاحظ، مثلًا، أنه إذا كانت $f(x,y,z)$ دالة
عودية جزئية ذات $3$ حجج، فإن الدالة
$g(x)\simeq f((x)_0,(x)_1,(x)_2)$ دالة عودية جزئية أحادية الحجة.
\end{explain}

\end{document}
